\documentclass[10pt,a4paper]{amsart}
\usepackage[margin=19mm]{geometry}
\usepackage{amsmath,amssymb,mathtools}
\usepackage[scr=rsfs,cal=euler]{mathalfa}
\usepackage{xfrac}
\usepackage[unicode,hidelinks,bookmarks=false]{hyperref}
\newcommand{\ie}{i.e.\ }
\newcommand{\cf}{cf.\ }
\newcommand{\resp}{respectively\ }
\newcommand{\Subsection}{\subsection}
\providecommand{\nobreakdash}{\nobreak\hbox{-}}
\setlength{\emergencystretch}{2em}
\multlinegap0pt
\usepackage{bm}
\usepackage{bbm}
\usepackage{dsfont}
\usepackage{xspace}
\usepackage{cancel}
\usepackage{mathtools}
\usepackage{amsthm}
\makeatletter
\@ifundefined{theorem}{\newtheorem{theorem}{Theorem}}{}
\makeatother
\def\ni{\noindent}
\title[Neighborhood Balanced k-Coloring of Graphs]{Neighborhood Balanced k-Coloring of Graphs}
\author{Maurice Genevieva Almeida, Tarkeshwar Singh, Siddharth Gupta, Ravindra Pawar}
\date{Source: arXiv:2509.06003v2 (CC BY 4.0)}
\begin{document}
\maketitle

\noindent\emph{Source and licence.} Neighborhood Balanced k-Coloring of Graphs. Version arXiv:2509.06003v2 (CC BY 4.0), distributed under the Creative Commons Attribution 4.0 International licence, as recorded in the arXiv licence metadata for this exact version and shown on the abstract page https://arxiv.org/abs/2509.06003v2. Full rights notice: licenses/arxiv-2509.06003-CC-BY-4.0.txt.\par\smallskip

\noindent\textit{Editorial scope.} Counted unit: the Theorem whose source label is \texttt{None} in the article (source0.tex lines 647--648, source label \texttt{None}), delivered together with the article's own complete original proof of it (source0.tex lines 650--718, one \texttt{proof} environment of 69 lines). No mathematics is written, repaired or re-derived: statement and proof are the article's own text line for line. Required context retained uncounted: none. Every definition, display and label of those lines is retained unchanged. Omitted from this package and not redistributed: 1--646, 649--649, 719--1545. Nothing inside a retained excerpt is omitted or altered: every retained body is delivered exactly as the article's lines are written, so no omission receipt of any kind accompanies this package. Citations: the retained excerpts carry no citation command at all, so no bibliography entry of the article is transcribed and no reference is redistributed. Reference adaptation: none was needed and none was made; every reference inside the delivered unit resolves inside it, so no unresolved reference or citation is produced. Printed slips kept literal: no printed word, symbol, hypothesis or number of the counted statement or of its proof was altered. Layout and class: the article's own class is \texttt{article} with packages including geometry, inputenc, babel, enumitem, ifthen, subcaption, lineno, xcolor, hyperref, fontawesome5, xparse, etoolbox, authblk, xspace; the wrapper is the lane's pinned \texttt{amsart} editorial wrapper at 10pt on A4, which restates the theorem-like environments the retained text needs and adds the article's own macro definitions that the retained text needs (\texttt{\textbackslash ni}), with extra wrapper packages (bm, bbm, dsfont, xspace, cancel, mathtools). The delivered numbering is the unit's own. Assets: no retained span contains a figure environment, a graphics-inclusion command, a PDF-page inclusion, a table or a TikZ picture; no image file is copied into this package, the package declares no asset dependency and no image file is redistributed with it. Licence: version arXiv:2509.06003v2 of this article is distributed under the Creative Commons Attribution 4.0 International licence, as recorded in the arXiv licence metadata for this exact version and shown on the abstract page https://arxiv.org/abs/2509.06003v2. A full rights notice is delivered at licenses/arxiv-2509.06003-CC-BY-4.0.txt. Characters: every retained excerpt is pure ASCII, so the delivered engine, XeLaTeX with its default Latin Modern text font, reports no missing character.
\begin{theorem} Let $G_1$ and $G_2$ be graphs admitting neighborhood-balanced $k_1$-coloring and neighborhood-balanced $k_2$-coloring, respectively. Then the direct product $G_1\times G_2$ admits a neighborhood-balanced $k_1k_2$-coloring.
\end{theorem}

\begin{proof}
Let
$f:V(G_1)\rightarrow \{1,2,\ldots,k_1\}$
be a neighborhood-balanced $k_1$-coloring of $G_1$, and let
$g:V(G_2)\rightarrow \{1,2,\ldots,k_2\}$
be a neighborhood-balanced $k_2$-coloring of $G_2$.
For each vertex $v\in V(G_2)$, consider the subgraph
\[
G_v=\{(u,v):u\in V(G_1)\},
\]
called the \emph{$G_1$-layer} corresponding to $v$.

\ni Suppose $g(v)=i$, where $1\le i\le k_2$. We color the vertices of the layer $G_v$ according to the coloring $f$ by assigning
\[
h(u,v)=k_1(i-1)+f(u).
\]
Thus, each $G_1$-layer corresponding to a vertex of color $i$ in $G_2$ is colored using the block of colors
\[
\{k_1(i-1)+1,\;k_1(i-1)+2,\;\ldots,\;k_1i\},
\]
following exactly the color pattern of the neighborhood-balanced coloring $f$ of $G_1$.
 Since the colors assigned to different layers of $G_1$ belong to disjoint blocks, the coloring $h$ uses exactly $k_1k_2$ distinct colors.\\
Now let $(x,y)\in V(G_1\times G_2)$, where $g(y)=i$. Since
\[
N_{G_1\times G_2}(x,y)
=
\{(u,v):u\in N_{G_1}(x),\,v\in N_{G_2}(y)\},
\]
every neighbor of $(x,y)$ belongs to a layer corresponding to a neighbor of $y$ in $G_2$.\\
\ni 
Fix an arbitrary color
$
c=k_1(j-1)+r$,
where $1\le j\le k_2$ and $1\le r\le k_1$.
A neighbor of $(x,y)$ has color $c$ if and only if its second coordinate is a neighbor of $y$ colored $j$ in $G_2$, and within the corresponding $G_1$-layer, its first coordinate is a neighbor of $x$ colored $r$ under $f$.

\ni
Since $g$ is a neighborhood-balanced $k_2$-coloring,
\[
\left|N_{G_2}(y)\cap g^{-1}(j)\right|
=
\frac{d_{G_2}(y)}{k_2}.
\]
Similarly, since $f$ is a neighborhood-balanced $k_1$-coloring,
\[
\left|N_{G_1}(x)\cap f^{-1}(r)\right|
=
\frac{d_{G_1}(x)}{k_1}.
\]
\ni
Therefore, the number of neighbors of $(x,y)$ having color $c$ is
\[
\frac{d_{G_2}(y)}{k_2}
\cdot
\frac{d_{G_1}(x)}{k_1}
=
\frac{d_{G_1}(x)d_{G_2}(y)}{k_1k_2}.
\]
\ni
Since
$d_{G_1\times G_2}(x,y)
=
d_{G_1}(x)\,d_{G_2}(y)$,
it follows that every color appears exactly
$\frac{d_{G_1\times G_2}(x,y)}{k_1k_2}$
times in the neighborhood of $(x,y)$.

\ni Hence, every vertex of $G_1\times G_2$ has an equal number of neighbors in each of the $k_1k_2$ colors. Therefore, $h$ is a neighborhood-balanced $k_1k_2$-coloring of $G_1\times G_2$.
\end{proof}

\end{document}
